
\section{Method}

\subsection{Overview}

Routing Attention consists of four steps:

\begin{lstlisting}
    Input X $\in$ R^{B$\times$L$\times$d}
            |
             
    +-------------------------+
    |  Q/K/V projections (LoRA)|
    +-------------+-----------+
                  |
                   
    +-------------------------+
    |  Compute joint feature  |
    |  F = normalize(Q + K)   |
    +-------------+-----------+
                  |
                   
    +-------------------------+
    |  k-means $\to$ cluster IDs  |
    +-------------+-----------+
                  |
                   
    +-------------------------+
    |  For each cluster:      |
    |    Mask to cluster + T0 |
    |    Softmax + weighted V |
    +-------------+-----------+
                  |
                   
    +-------------------------+
    |  Output projection      |
    +-------------------------+
\end{lstlisting}

\textbf{Key idea.} Each query attends only to keys in its cluster (plus the first 
token). Clusters are computed per-input via k-means.

\subsection{Cluster Assignment}

\subsubsection{Joint Feature}

We compute a joint representation of queries and keys:

\[
F = \text{normalize}(Q + K) \in \mathbb{R}^{B \times H \times L \times d}
\tag{4.1}
\]

where \texttt{normalize} is $\ell_2$ normalization along the head dimension. The 
$Q + K$ sum captures both "who is asking" and "who is being asked", 
providing a richer feature than Q alone.

\textbf{GQA handling.} If $H_{kv} < H$, we repeat K to match H before summing.

\subsubsection{Online k-Means}

We use Lloyd's algorithm for $n_{\text{iters}}$ iterations:

\textbf{Algorithm 1:} Cluster assignment
\begin{lstlisting}
Input:  F $\in$ R^{B$\times$H$\times$L$\times$d}, K clusters, n\_iters
Output: cluster\_ids $\in$ {0, ..., K-1}^{B$\times$H$\times$L}

1: step $\leftarrow$ max(1, L // K)
2: init\_idx $\leftarrow$ [0, step, 2*step, ..., (K-1)*step]    Deterministic init
3: centroids $\leftarrow$ F[:, :, init\_idx, :]                  (B, H, K, d)
4:
5: for iter = 1 to n\_iters do
6:     dist $\leftarrow$ cdist(F, centroids)                     (B, H, L, K)
7:     cluster\_ids $\leftarrow$ argmin(dist, dim=-1)
8:     one\_hot $\leftarrow$ one\_hot(cluster\_ids, K)
9:     counts $\leftarrow$ one\_hot.sum(dim=2).clamp(min=1)       (B, H, K)
10:    sum\_f $\leftarrow$ einsum("bhlk,bhld-\textgreater{}bhkd", one\_hot, F)
11:    centroids $\leftarrow$ sum\_f / counts
12: end for
13: return cluster\_ids
\end{lstlisting}

\textbf{Initialization.} We use deterministic stride-based initialization (every 
\texttt{step}-th token). This avoids the randomness of k-means++ and makes 
clustering reproducible.

\textbf{Edge case.} If $L < K$, we reduce $K$ to $L$.

\subsection{Cluster-Restricted Attention}

For each cluster $k$, we compute attention over its tokens:

\[
\mathcal{S}\_k = \{j : c\_j = k\} \cup \{0\}
\tag{4.2}
\]

The union with $\{0\}$ adds the first token as a global context anchor.

The attention output for cluster $k$:

\[
\text{out}\_i^{(k)} = \text{softmax}\!\left(\frac{q\_i K\_{\mathcal{S}\_k}^\top}{\sqrt{d}}\right) V\_{\mathcal{S}\_k} \quad \text{if } i \in \mathcal{S}\_k
\tag{4.3}
\]

\subsubsection{Masking Implementation}

In practice, we compute the \textbf{full} $L \times L$ attention matrix and mask 
it:

\begin{lstlisting}
for k in range(K\_num):
    mask\_k = (cluster\_ids == k) | (torch.arange(L).view(1, 1, -1) == 0)
    mask\_k = mask\_k & valid\_mask\_2d.unsqueeze(1)
    
    scores = torch.matmul(Q, K.transpose(-2, -1)) * scale
    scores = scores.masked\_fill(~mask\_k.unsqueeze(2), -1e4)
    
    # Self-token preservation
    eye = torch.eye(L, dtype=torch.bool)
    scores = scores.masked\_fill(eye, 0.0)
    
    if is\_causal:
        causal\_mask = torch.triu(torch.full((L, L), True), diagonal=1)
        scores = scores.masked\_fill(causal\_mask & (~eye), -1e4)
    
    attn = F.softmax(scores, dim=-1)
    out\_k = torch.matmul(attn, V)
    outputs = torch.where(mask\_k.unsqueeze(-1), out\_k, outputs)
\end{lstlisting}

\textbf{Cost.} This computes $K$ full attention matrices. \textbf{Slower} than dense.

\subsubsection{Self-Token Preservation}

We force each token to attend to itself by setting the diagonal of the score 
matrix to 0. This prevents the case where a token is isolated from its own 
attention.

\subsubsection{Causal Masking}

For decoders, we apply an upper-triangular mask (standard causal):

\[
\text{causal\_mask}\_{ij} = \begin{cases} 0 & j \leq i \\ -\infty & j \textgreater{} i \end{cases}
\tag{4.4}
\]

Combined with the cluster mask:

\[
\text{final\_mask} = \text{cluster\_mask} \wedge \text{causal\_mask}
\tag{4.5}
\]

\subsubsection{Token-0 Anchor}

Token 0 is added to every cluster's key set. This acts as a **universal 
context anchor**, since the first token often accumulates summary 
information (attention sink phenomenon, Xiao et al., 2024).

\textbf{Why is this necessary?} Without it, some clusters may contain only 
"uninformative" tokens, and the attention output may be dominated by 
noise. Token 0 provides a stable reference.

\subsection{Output Combination}

The final output is a masked sum over clusters:

\[
\text{out}\_i = \sum\_k \mathbb{1}[c\_i = k] \cdot \text{out}\_i^{(k)}
\tag{4.6}
\]

Only the cluster containing token $i$ contributes. This is implemented as 
\texttt{torch.where(mask\_k.unsqueeze(-1), out\_k, outputs)}.

\subsection{Padding Masking}

Padding tokens do not contribute to any cluster's attention:

\begin{lstlisting}
mask\_k = mask\_k & valid\_mask\_2d.unsqueeze(1)
\end{lstlisting}

Padding positions in the output are zeroed.

\subsection{LoRA Adaptation}

Each projection is a \texttt{LoRALinear}:

\begin{lstlisting}
out = F.linear(x, W\_frozen, b) + F.linear(F.linear(x, A), B) * scaling
\end{lstlisting}

Only \texttt{A, B} are trainable; \texttt{W\_frozen} is copied from the pretrained attention.

\subsection{Algorithm}

\textbf{Algorithm 2:} Routing Attention forward pass
\begin{lstlisting}
Input:  X $\in$ R^{B$\times$L$\times$d}, attention\_mask
Output: Y $\in$ R^{B$\times$L$\times$d}

1: Q $\leftarrow$ X W\_Q;  K $\leftarrow$ X W\_K;  V $\leftarrow$ X W\_V        LoRA projections
2: if RoPE: (Q, K) $\leftarrow$ Rotary(Q, K)
3: if GQA:  (K, V) $\leftarrow$ repeat\_kv(K, V)
4:
5: F $\leftarrow$ normalize(Q + K)
6: cluster\_ids $\leftarrow$ kmeans(F, K, n\_iters)
7:
8: outputs $\leftarrow$ zeros\_like(Q)
9: for k = 0 to K-1 do
10:     mask\_k $\leftarrow$ (cluster\_ids == k) | (arange(L) == 0)
11:     mask\_k $\leftarrow$ mask\_k & valid\_mask
12:
13:     scores $\leftarrow$ Q K^T / $\sqrt{}$d
14:     scores $\leftarrow$ masked\_fill(scores, ~mask\_k, -1e4)
15:     scores $\leftarrow$ masked\_fill(scores, eye, 0)        Self-token
16:     if causal: scores $\leftarrow$ masked\_fill(scores, triu & ~eye, -1e4)
17:     attn $\leftarrow$ softmax(scores)
18:     out\_k $\leftarrow$ attn V
19:     outputs $\leftarrow$ where(mask\_k, out\_k, outputs)
20: end for
21:
22: if not is\_nested\_encoder: outputs $\leftarrow$ out\_proj(outputs)
23: return outputs
\end{lstlisting}

\subsection{Complexity Analysis}

\begin{table}[H]
\centering
\small
\begin{tabular}{lll}
\toprule
\textbf{Component} & \textbf{Time} & \textbf{Memory} \\
\midrule
Q/K/V projection & $\mathcal{O}(BLd^2)$ & $\mathcal{O}(BLd)$ \\
k-means (per iter) & $\mathcal{O}(BHKLd)$ & $\mathcal{O}(BHKL)$ \\
Attention loop & $\mathcal{O}(K \cdot BHL^2)$ & $\mathcal{O}(BHL^2)$ \\
Output projection & $\mathcal{O}(BLd^2)$ & $\mathcal{O}(BLd)$ \\
\bottomrule
\end{tabular}
\end{table}

\textbf{The attention loop is $\mathcal{O}(K \cdot BHL^2)$} -- $K$ times slower 
than dense attention.

\textbf{Theoretical sparse cost} (with sparse matmul) would be $\mathcal{O}(L^2/K)$ 
per head. Our implementation does not use sparse matmul; it computes the 
full matrix per cluster.

\subsection{Parameter Analysis}

Routing Attention adds LoRA adapters (no new architectural parameters):

\begin{table}[H]
\centering
\small
\begin{tabular}{lll}
\toprule
\textbf{Component} & \textbf{Per layer} & \textbf{Total (DistilBERT)} \\
\midrule
LoRA (4 proj $\times$ $r$) & 4 $\times$ 2 r d & 6 $\times$ 98,304 = 589,824 \\
\textbf{Total added} &  & \textbf{~589,824} \\
\textbf{Percentage} &  & \textbf{+0.88\%} \\
\bottomrule
\end{tabular}
\end{table}

\textbf{Routing Attention increases parameter count.} It is not a compressor.

\subsection{Design Choices}

\subsubsection{Why Q+K Instead of Q?}

Using Q alone (as in Reformer or Routing Transformer) ignores key content. 
Query and key are both learned projections of the same input, so their sum 
captures more information.

\subsubsection{Why Deterministic Init?}

k-means++ introduces randomness. Deterministic stride init makes runs 
reproducible without loss of quality (in our experiments).

\subsubsection{Why Token-0 Anchor?}

Without an anchor, clusters with sparse or noisy content may produce 
degenerate attention outputs. Token 0 provides a stable reference.

\subsubsection{Why Self-Token Preservation?}

Forcing $j = i$ in the attention pattern prevents a query from being 
disconnected from its own representation.
\begin{lstlisting}

---
\end{lstlisting}